% English narration and Vietnamese production notes, maintained source.
\documentclass[10pt,a4paper]{article}
\usepackage{fontspec}
\IfFontExistsTF{TeX Gyre Heros}{\setmainfont{TeX Gyre Heros}}{\setmainfont{Arial}}
\usepackage[margin=20mm]{geometry}
\usepackage{xcolor,hyperref,needspace}
\definecolor{navy}{RGB}{18,30,49}
\hypersetup{colorlinks=true,urlcolor=navy,pdftitle={Divergency - English Voice Script},pdfauthor={Divergency Team}}
\setlength{\parindent}{0pt}
\setlength{\parskip}{3pt}
\setlength{\emergencystretch}{3em}
\begin{document}

Kịch bản devlog Divergency\par

English voice script with Vietnamese production notes 29/09/2026 | Video 10–12 phút | Mục tiêu dựng 11:30 Tài liệu dùng để thu lời kể, gom tư liệu và dựng video về hành trình hình thành Divergency. Mạch chính đi từ cảm hứng LF2 đến quá trình biến một nhân vật thành đòn đánh chơi được, rồi dẫn tới bước chuẩn bị Kickstarter.\par

\textbf{Cách dùng:}  Toàn bộ phần English voiceover là lời tiếng Anh để thu. Hướng dẫn dựng giữ tiếng Việt. Bản tiếng Việt gốc được giữ riêng. Các phần Hình, Dựng, Âm thanh và Chữ là chỉ dẫn sản xuất. Mã F01 đến F23 dẫn tới checklist tư liệu. Mốc thời gian là dự kiến; căn lại theo bản thu thật.\par

\textbf{Danh tính dự án:}  Divergency là game riêng của nhóm. Little Fighter 2 chỉ là nguồn cảm hứng được ghi tên trong phần tư liệu tham khảo.\par

Mục lục\par

\Needspace{3\baselineskip}\textbf{1 Nhận xét và hướng biên tập 2}\par

\Needspace{3\baselineskip}\textbf{2 Mạch video 2}\par

\Needspace{3\baselineskip}\textbf{3 Một đòn đánh, sáu năm 3}\par

\Needspace{3\baselineskip}\textbf{4 Nguồn cảm hứng và thế giới Divergency 4}\par

\Needspace{3\baselineskip}\textbf{5 Một dự án nằm giữa cuộc sống 5}\par

\Needspace{3\baselineskip}\textbf{6 Tại sao lại là Divergency? 6}\par

\Needspace{3\baselineskip}\textbf{7 Một thế giới tối, một người kể chuyện vui 7}\par

\Needspace{3\baselineskip}\textbf{8 Từ hình vẽ đến chuyển động 8}\par

\Needspace{3\baselineskip}\textbf{9 Đẹp thôi chưa đủ: phải đánh được 9}\par

\Needspace{3\baselineskip}\textbf{10 Nhảy lên dễ, đáp xuống mới khó 10}\par

\Needspace{3\baselineskip}\textbf{11 Cuối cùng, một đòn đánh ra hồn 11}\par

\Needspace{3\baselineskip}\textbf{12 Sáu năm có nghĩa là gì? 12}\par

\Needspace{3\baselineskip}\textbf{13 Bước tiếp theo: Kickstarter 14}\par

\Needspace{3\baselineskip}\textbf{14 Checklist thu thập tư liệu 14}\par

\Needspace{3\baselineskip}\textbf{15 Hướng dẫn dựng 18}\par

\Needspace{3\baselineskip}\textbf{16 Thứ tự thu thập và quay 19}\par

\Needspace{3\baselineskip}\textbf{17 Lời mời Kickstarter theo trạng thái trang 19}\par

\Needspace{3\baselineskip}\textbf{18 Chi tiết cần khớp trước khi xuất 19}\par

\Needspace{3\baselineskip}\textbf{19 Điểm tìm tư liệu trong thư mục dự án 20}\par

\Needspace{3\baselineskip}\textbf{20 Tiêu đề và thumbnail 20}\par

\Needspace{3\baselineskip}\textbf{21 Kiểm tra trước khi xuất 21}\par

\clearpage

\Needspace{3\baselineskip}\textbf{1 Nhận xét và hướng biên tập}\par

\textbf{Điểm nên giữ:}  Nguồn cảm hứng LF2, sự hợp tác với Khoa, quãng sáu năm và cái tên gắn với công việc nghiên cứu. Đây là những phần khiến câu chuyện có tác giả cụ thể.\par

\textbf{Mở bằng một kết quả:}  Đưa game hiện tại và bản cũ vào 12 giây đầu. Người xem biết mình sắp nghe câu chuyện về thứ gì trước khi nghe tên công cụ.\par

\textbf{Tổ chức lại phần kỹ thuật:}  Theo một nhân vật và một đòn đánh: hình vẽ → chuyển động → vùng đánh → tiếp đất → gameplay. Mỗi phần trả lời một câu hỏi nhìn thấy được.\par

\textbf{Giữ sự tự tin khi kể:}  Dùng “At the time, I wasn’t confident building the whole tool myself” thay cho tự hạ thấp năng lực. Trải nghiệm khó vẫn thật, nhưng câu chuyện cho thấy bạn tìm được cách làm.\par

\textbf{Kết bằng điều có thể thấy:}  Cho người xem nhìn cặp ảnh cũ/mới và gameplay trước khi nói tới tự hào và Kickstarter. Lời mời ủng hộ khi đó có cơ sở cụ thể.\par

\Needspace{3\baselineskip}\textbf{2 Mạch video}\par

Thời gian Nội dung\par

00:00–00:35 Một đòn đánh, sáu năm\par

00:35–01:50 Nguồn cảm hứng và thế giới Divergency\par

01:50–02:35 Một dự án nằm giữa cuộc sống\par

02:35–03:35 T ại sao lại là Divergency?\par

03:35–04:30 Một thế giới tối, một người kể chuyện vui\par

04:30–05:45 Từ hình vẽ đến chuyển động\par

05:45–07:10 Đẹp thôi chưa đủ: phải đánh được\par

07:10–08:05 Nhảy lên dễ, đáp xuống mới khó\par

08:05–09:10 Cuối cùng, một đòn đánh ra hồn\par

09:10–10:35 Sáu năm có nghĩa là gì?\par

10:35–11:30 Bước tiếp theo: Kickstarter\par

\clearpage

\Needspace{3\baselineskip}\textbf{3 Một đòn đánh, sáu năm}\par

00:00–00:35 — Cho xem điều đáng quan tâm trước khi kể lịch sử.\par

\Needspace{3\baselineskip}\textbf{S01 Game trước, lời giới thiệu sau}\par

00:00–00:12 · 12 giây · Tư liệu F01, F02\par

\textbf{English voiceover:}  This is Divergency, the action game a friend and I are making. And this... is what one of our early attempts looked like.\par

\textbf{Hình cần lấy:}  0–5s: một pha đánh rõ, đẹp từ build hiện tại. 5–9s: cắt sang thử nghiệm cũ tương\par

\textbf{phản rõ. 9–12s:}  giữ một lỗi nhỏ thật hoặc pose vụng đáng yêu.\par

\textbf{Dựng:}  Cắt theo điểm chạm của đòn đánh; hai bản cùng vị trí nhân vật càng tốt. Không intro logo dài.\par

\textbf{Âm thanh:}  Mở bằng tiếng game 1–2s; nhạc có nhịp nhẹ vào dưới giọng. Hạ nhạc ở dấu ngập ngừng.\par

\textbf{Chữ trên hình:}  DIVERGENCY · work in progress → Early prototype\par

\textbf{Ghi chú:}  Chỉ dùng lỗi có tư liệu. Nếu không có, so một pose cũ với pose mới.\par

\Needspace{3\baselineskip}\textbf{S02 Sáu năm không nằm trong một cú click}\par

00:12–00:23 · 11 giây · Tư liệu F02, F06\par

\textbf{English voiceover:}  There’s about six years between those clips. Sometimes we made a little progress. Sometimes the project sat untouched for a year. But we kept coming back.\par

\textbf{Hình cần lấy:}  Ảnh cũ → màn hình project hiện tại → tay bấm Play. Một lát cắt bàn làm việc, không cần dựng montage lịch dày đặc.\par

\textbf{Dựng:}  Dùng cùng một khung màn hình để nối cũ và mới. Giữ hình cuối thêm một nhịp.\par

\textbf{Âm thanh:}  Nhạc ấm lên, không chuyển sang piano bi lụy.\par

\textbf{Chữ trên hình:}  6 years · many fresh starts\par

\Needspace{3\baselineskip}\textbf{S03 Lời hứa với người xem}\par

00:23–00:35 · 12 giây · Tư liệu F01, F16, F21, F22\par

\textbf{English voiceover:}  Today, I want to share how we started building Divergency and developed its own identity. And what it took to make our characters fight... and actually stay on the ground.\par

\textbf{Hình cần lấy:}  Cận mặt kể chuyện hoặc tay bên bàn phím → 2s nhân vật đánh → 2s tiếp đất → logo nhỏ trên khung game.\par

\textbf{Dựng:}  Đừng cho xem hết cú lỗi tiếp đất; để dành phần giải thích. Logo xuất hiện cùng hình, khoảng 1s.\par

\textbf{Âm thanh:}  Cười nhẹ ở câu cuối; tiếng click chuyển sang phần nguồn cảm hứng.\par

\textbf{Chữ trên hình:}  From an idea to a working attack\par

\clearpage

\Needspace{3\baselineskip}\textbf{4 Nguồn cảm hứng và thế giới Divergency}\par

00:35–01:50 — Hoài niệm → thử nghiệm → cùng Khoa chọn hướng riêng.\par

\Needspace{3\baselineskip}\textbf{S04 Tình yêu bắt đầu từ LF2}\par

00:35–00:50 · 15 giây · Tư liệu F03, F06\par

\textbf{English voiceover:}  It started with Little Fighter 2. I loved how much action could fit on one little screen. Eventually, I had the thought many fans probably have: what if I tried making something myself?\par

\textbf{Hình cần lấy:}  Gameplay LF2 tự quay: chọn một tình huống dễ đọc, không cần montage tên tất cả nhân vật. Có thể chèn 2s tay chơi trên bàn phím.\par

\textbf{Dựng:}  Lời dẫn vào trước hình LF2 một chút để chuyển tự nhiên. Ghi tên game rõ.\par

\textbf{Âm thanh:}  Tiếng phím/tiếng game thấp dưới giọng, nhạc vui nhưng bớt dồn.\par

\textbf{Chữ trên hình:}  Inspired by Little Fighter 2\par

\Needspace{3\baselineskip}\textbf{S05 Các lần thử đầu tiên}\par

00:50–01:05 · 15 giây · Tư liệu F02\par

\textbf{English voiceover:}  At first, I just experimented. Make a character run. Try an attack. Build it and see what happens. Looking back, it’s pretty rough. But getting that first character to work felt amazing.\par

\textbf{Hình cần lấy:}  Ba clip cũ khoảng 4s mỗi clip: di chuyển, đánh, bản build. 3s cuối giữ chi tiết tác giả tự hào.\par

\textbf{Dựng:}  Đặt tên phiên bản hoặc năm khi có bằng chứng. Giữ tỉ lệ nguồn, dùng nền đơn giản nếu video cũ nhỏ.\par

\textbf{Âm thanh:}  Giữ một tiếng game cũ cho cảm giác tư liệu; không phủ meme sound.\par

\textbf{Chữ trên hình:}  Our early prototypes\par

\Needspace{3\baselineskip}\textbf{S06 Rủ Khoa cùng làm}\par

01:05–01:20 · 15 giây · Tư liệu F04, F08\par

\textbf{English voiceover:}  Then I started talking with Khoa. A few ideas about characters and combat slowly grew into something bigger. We wanted our own world, with its own atmosphere and a reason for these characters to exist.\par

\textbf{Hình cần lấy:}  Hai bong bóng hội thoại, mỗi lần chỉ một câu. Lời tiếng Anh đề xuất: “What if we built our own world?” → “Sure. But let’s make it a little darker.” Sau đó mở concept.\par

\textbf{Dựng:}  Bong bóng trượt nhẹ 6–10 frame rồi dừng; đừng mô phỏng màn hình chat dài. Nhãn dựng lại đặt ngay khi xuất hiện.\par

\textbf{Âm thanh:}  Hai tiếng pop nhỏ, không dùng giọng giả Khoa. Người kể tiếp tục bằng giọng mình.\par

\textbf{Chữ trên hình:}  Recreated conversation · rewritten dialogue\par

\textbf{Ghi chú:}  Đây là lời thoại đề xuất, không phải trích chat thật. Nếu dùng chat thật, lấy đúng câu và ngữ cảnh.\par

\Needspace{3\baselineskip}\textbf{S07 Từ cảm hứng tới Divergency}\par

01:20–01:35 · 15 giây · Tư liệu F02, F08, F09, F21\par

\textbf{English voiceover:}  That’s how Divergency started taking shape. As we developed the art, the world, and the gameplay, we kept defining its own identity.\par

\textbf{Hình cần lấy:}  Một thử nghiệm cũ → concept thế giới → Deep/Solei → logo. Mỗi hình hỗ trợ đúng danh từ đang nói.\par

\textbf{Dựng:}  Bong bóng cuối thu về góc, concept mở ra từ cùng vị trí. Ưu tiên cắt đơn giản nếu thao tác này tốn công.\par

\textbf{Âm thanh:}  Nhạc tiến thêm một lớp, chưa đạt cao trào.\par

\textbf{Chữ trên hình:}  A world of our own\par

\clearpage

\Needspace{3\baselineskip}\textbf{S08 Một lỗi nhỏ, một bước tiến của Divergency}\par

01:35–01:50 · 15 giây · Tư liệu F05, F22, F06\par

\textbf{English voiceover:}  Sometimes I fixed one bug, ran the game again, and found another. Each time, I understood Divergency’s characters and combat a little better. My development speed? That’s another story.\par

\textbf{Hình cần lấy:}  Hai clip thật trước và sau khi sửa cùng một lỗi trong Divergency → gương mặt cười chừ → build hiện tại.\par

\textbf{Dựng:}  Một nhịp dừng trước câu cuối. Dùng clip lỗi và clip sau khi sửa của cùng một cảnh Divergency.\par

\textbf{Âm thanh:}  Hạ nhạc ngắn ở câu đùa, rồi nối bằng tiếng gõ phím.\par

\textbf{Chữ trên hình:}  One bug fixed · one step forward\par

\textbf{Ghi chú:}  Dùng lỗi đã quay thật của Divergency; không trình bày cảnh dựng lại như lịch sử phát triển.\par

\Needspace{3\baselineskip}\textbf{5 Một dự án nằm giữa cuộc sống}\par

01:50–02:35 — Thành thật về quãng ngắt, rồi trở lại hiện tại.\par

\Needspace{3\baselineskip}\textbf{S09 Làm game sau giờ làm}\par

01:50–02:05 · 15 giây · Tư liệu F06\par

\textbf{English voiceover:}  Because we still had our day jobs. The game had to fit around them: an evening, a weekend, a few hours when we had the energy. Some days, we barely moved forward.\par

\textbf{Hình cần lấy:}  Cảnh trở về bàn làm việc → bật đèn → mở project. Quay ba góc trong cùng buổi, không cần đóng cảnh kiệt sức.\par

\textbf{Dựng:}  Tiếng phím từ chương trước giữ qua điểm cắt. Khung hình dịu và tĩnh hơn 15–20s.\par

\textbf{Âm thanh:}  Nhạc bớt bộ gõ, giữ chất ấm. Room tone nhẹ.\par

\textbf{Chữ trên hình:}  After work\par

\Needspace{3\baselineskip}\textbf{S10 Khoảng ngắt là một phần câu chuyện}\par

02:05–02:20 · 15 giây · Tư liệu F02, F06\par

\textbf{English voiceover:}  There were stretches when we barely touched the game for a whole year. Coming back meant remembering what I was even doing. The ideas were still there. Finding the time was harder.\par

\textbf{Hình cần lấy:}  Ảnh project cũ hoặc tài liệu thật → tay cuộn lại ghi chú → mở scene hiện tại. Một cảnh giữ 5–6s cho người xem thở.\par

\textbf{Dựng:}  Dissolve ngắn chỉ khi đổi thời điểm. Không chạy lịch giả để chứng minh quãng nghỉ.\par

\textbf{Âm thanh:}  Giọng bình thản, có thể mỉm cười ở “what I was even doing”.\par

\textbf{Chữ trên hình:}  Some very long breaks\par

\Needspace{3\baselineskip}\textbf{S11 Điểm quay lại}\par

02:20–02:35 · 15 giây · Tư liệu F06, F04, F01\par

\textbf{English voiceover:}  But we’d open it again. Fix something small. Try an idea. Send each other a new drawing. I think that’s what kept it going: we still wanted to see this world come to life.\par

\textbf{Hình cần lấy:}  Một thay đổi nhỏ thực tế → gửi ảnh trong phần dựng lại hoặc xem ảnh chung → nhân vật chạy trong scene.\par

\textbf{Dựng:}  Cắt theo cú bấm Play; nhịp hình bắt đầu sáng và linh hoạt hơn.\par

\textbf{Âm thanh:}  Đưa bộ gõ nhẹ trở lại khi nói “open it again”.\par

\textbf{Chữ trên hình:}  Opening the project again\par

\clearpage

\Needspace{3\baselineskip}\textbf{6 Tại sao lại là Divergency?}\par

02:35–03:35 — Dùng nghề nghiên cứu làm một câu chuyện cá nhân dễ hiểu.\par

\Needspace{3\baselineskip}\textbf{S12 Công việc ban ngày của tác giả}\par

02:35–02:50 · 15 giây · Tư liệu F07, F06\par

\textbf{English voiceover:}  The name has a funny connection to my work. My PhD research is in control\par

\textbf{systems:}  controllers, observers, stability, convergence... basically, a lot of time trying to make systems behave.\par

\textbf{Hình cần lấy:}  Bàn nghiên cứu hoặc màn hình mô phỏng → lần lượt hiện ba nhãn, không đưa nguyên trang phương trình.\par

\textbf{Dựng:}  Giọng vẫn kể chuyện, đừng đổi thành thuyết trình. Chỉ một sơ đồ ở trung tâm.\par

\textbf{Âm thanh:}  Nhạc có nhịp gõ nhỏ đều, gợi cảm giác sắp xếp.\par

\textbf{Chữ trên hình:}  Control systems → stability → convergence\par

\textbf{Ghi chú:}  Thu đúng tư cách hiện tại. Nếu câu về PhD không phù hợp, dùng: “I work in control systems\par

\textbf{research:}  controllers, observers, stability, convergence... basically, a lot of time trying to make systems behave.”\par

\Needspace{3\baselineskip}\textbf{S13 Hội tụ, nói theo cách dễ hình dung}\par

02:50–03:05 · 15 giây · Tư liệu F07, F23\par

\textbf{English voiceover:}  For a problem like this, I want the system to reach a desired state. It looks wonderfully organized in a diagram. In my simulations, though... sometimes the curve has other plans.\par

\textbf{Hình cần lấy:}  Đồ thị minh họa hội tụ về mục tiêu, rồi một đường lệch ra khỏi khung. Nếu có simulation thật phù hợp thì dùng thay.\par

\textbf{Dựng:}  Đường đi ra khỏi khung làm cú đùa thị giác duy nhất; tránh rung cả màn hình.\par

\textbf{Âm thanh:}  Một nốt trượt nhẹ hoặc tiếng click lỗi nhỏ; chừa nửa nhịp sau câu cuối.\par

\textbf{Chữ trên hình:}  Illustration · not experimental data\par

\Needspace{3\baselineskip}\textbf{S14 Một bên muốn ổn định, một bên thích bất ổn}\par

03:05–03:20 · 15 giây · Tư liệu F07, F08, F23\par

\textbf{English voiceover:}  Meanwhile, Khoa wanted a darker, more unstable world, where things could drift away from the path you’d expect. I spent my days looking for stability. Our game was heading somewhere else.\par

\textbf{Hình cần lấy:}  Nửa trái là đồ thị gọn; nửa phải là concept. Hai khung nhập thành một artwork toàn màn hình.\par

\textbf{Dựng:}  Đường đồ thị nối với đường chân trời/đường sáng có sẵn trong artwork nếu hợp bố cục. Không ép vẽ đè lên art.\par

\textbf{Âm thanh:}  Thêm một lớp ambience của thế giới game dưới nhạc.\par

\textbf{Chữ trên hình:}  Me: looking for stability · Khoa: adding instability\par

\Needspace{3\baselineskip}\textbf{S15 Tên game tìm được chỗ đứng}\par

03:20–03:35 · 15 giây · Tư liệu F21, F08, F23\par

\textbf{English voiceover:}  That’s the connection behind Divergency: my research meeting the world we wanted to build. During the day, I work on convergence. At night, I work on... Divergency.\par

\textbf{Hình cần lấy:}  Chữ Convergence tách nhẹ thành các nét → tên Divergency hiện đầy đủ → cảnh môi trường giữ 4–5s.\par

\textbf{Dựng:}  Đây là ẩn dụ cá nhân, không trình bày thành định nghĩa toán. Logo không cần biến hình phức tạp.\par

\textbf{Âm thanh:}  Ngắt nhẹ trước tên game. Giữ 2s chỉ ambience cuối cảnh.\par

\textbf{Chữ trên hình:}  DIVERGENCY\par

\clearpage

\Needspace{3\baselineskip}\textbf{7 Một thế giới tối, một người kể chuyện vui}\par

03:35–04:30 — Cho thấy cá tính hình ảnh và hai nhân vật.\par

\Needspace{3\baselineskip}\textbf{S16 Khán giả được bước vào thế giới}\par

03:35–03:50 · 15 giây · Tư liệu F08\par

\textbf{English voiceover:}  We’ve developed concept art for that world. It’s dark, and some places feel unsettling before anything even happens. Let me show you a few pieces we’re really happy with.\par

\textbf{Hình cần lấy:}  Ba concept tuyển kỹ: toàn cảnh 5s → trung cảnh 5s → chi tiết 5s. Chọn hình khác bố cục, không chỉ đổi màu.\par

\textbf{Dựng:}  Zoom chậm rất nhẹ trên ảnh lớn; khi đổi hình giữ một điểm sáng gần cùng vị trí.\par

\textbf{Âm thanh:}  Giảm nhạc, cho ambience mỗi không gian nổi nhẹ ở khe lời.\par

\textbf{Chữ trên hình:}  Concept art · a world in development\par

\Needspace{3\baselineskip}\textbf{S17 Tối nhưng vẫn phải nhìn rõ}\par

03:50–04:05 · 15 giây · Tư liệu F08, F01\par

\textbf{English voiceover:}  There are some pretty brutal moments, some blood, and a few... unusual designs. But first, I want you to get a feel for the world and its personality. We’ll get to the surprises.\par

\textbf{Hình cần lấy:}  Toàn cảnh môi trường → silhouette địch → một đòn của Deep ở khung rộng. Nếu cần chi tiết máu, dùng ngắn và để hành động đọc được.\par

\textbf{Dựng:}  Không dùng máu hoặc cú zoom cơ thể làm dấu câu cho mọi lần đổi hình.\par

\textbf{Âm thanh:}  Tiếng đòn đánh đúng một lần; tránh tăng âm đột ngột.\par

\textbf{Chữ trên hình:}  Atmosphere · characters · movement\par

\Needspace{3\baselineskip}\textbf{S18 Một câu đùa hình ảnh tùy chọn}\par

04:05–04:17 · 12 giây · Tư liệu F08, F22\par

\textbf{English voiceover:}  And here’s a design that made me pause for a second. Well... I think you can see what I mean by surprises.\par

\textbf{Hình cần lấy:}  Nếu hợp mức độ thoải mái của bạn: boss áo trắng ở toàn thân 2–3s, phản ứng người kể 1s, rồi chuyển sang nhân vật chính. Phương án mặc định nhẹ: một silhouette kỳ quặc.\par

\textbf{Dựng:}  Chỉ một nhịp hài. Nếu chi tiết hở làm bạn thấy thô, bỏ hẳn và giữ 12s cho concept đẹp.\par

\textbf{Âm thanh:}  Ngưng lời nửa nhịp; không dùng tiếng huýt sáo hoặc hiệu ứng trêu cơ thể.\par

\textbf{Chữ trên hình:}  Không cần chữ thêm\par

\textbf{Ghi chú:}  Cảnh tùy chọn. Chưa tìm/xem tư liệu boss cụ thể; không mặc định asset nào là cảnh này.\par

\Needspace{3\baselineskip}\textbf{S19 Gặp Deep và Solei}\par

04:17–04:30 · 13 giây · Tư liệu F09, F10\par

\textbf{English voiceover:}  These are Deep and Solei, two of our main characters. Seeing the finished artwork felt great. Getting them off the canvas and into a fight... took a little more work.\par

\textbf{Hình cần lấy:}  Hai nhân vật trên nền sạch, tên rõ → chọn một nhân vật làm ví dụ xuyên suốt phần kỹ thuật → giữ một pose.\par

\textbf{Dựng:}  Cắt đúng pose từ artwork sang canvas của phần mềm ở cảnh sau. Một nhân vật xuyên suốt giúp người xem theo dễ.\par

\textbf{Âm thanh:}  Ambience nhường cho nhạc nhịp nhẹ và tiếng chuột.\par

\textbf{Chữ trên hình:}  DEEP · SOLEI\par

\clearpage

\Needspace{3\baselineskip}\textbf{8 Từ hình vẽ đến chuyển động}\par

04:30–05:45 — Photoshop → Aseprite → Unity, kể bằng trải nghiệm của nhóm.\par

\Needspace{3\baselineskip}\textbf{S20 Hồi còn làm bằng Photoshop}\par

04:30–04:45 · 15 giây · Tư liệu F11, F10\par

\textbf{English voiceover:}  We started in Photoshop. On my machine, with our workflow, the files got heavy and managing all those frames became tiring. Rearranging sprites could eat up time before I’d really changed the animation.\par

\textbf{Hình cần lấy:}  Project Photoshop thật hoặc ảnh giao diện cũ → các frame đang sắp xếp → một thao tác lặp lại. Chỉ cần cho xem một bất tiện cụ thể.\par

\textbf{Dựng:}  Không dựng thanh loading giả để chê phần mềm. Nếu không còn project cũ, kể trên sprite sheet hiện có.\par

\textbf{Âm thanh:}  Nhạc nhẹ, tiếng click thưa. Không kéo dài cảm giác bực bội.\par

\textbf{Chữ trên hình:}  Our earlier workflow\par

\Needspace{3\baselineskip}\textbf{S21 Gặp Aseprite}\par

04:45–05:00 · 15 giây · Tư liệu F12\par

\textbf{English voiceover:}  Then I found Aseprite. For our sprite work, it felt much more comfortable. I could adjust a frame and preview it quickly. That shorter loop made me want to keep working.\par

\textbf{Hình cần lấy:}  Mở Aseprite → chọn đúng animation đang kể → đổi qua vài frame → bấm Play. Quay con trỏ chậm, dễ theo.\par

\textbf{Dựng:}  Giữ nhân vật cùng vị trí với cảnh trước. Có thể đẩy gần timeline một lần rồi trở lại toàn khung.\par

\textbf{Âm thanh:}  Nhạc sáng hơn một chút. Một tiếng Play/click làm điểm mở.\par

\textbf{Chữ trên hình:}  ASEPRITE\par

\Needspace{3\baselineskip}\textbf{S22 Cho xem một thay đổi thật}\par

05:00–05:15 · 15 giây · Tư liệu F12, F10\par

\textbf{English voiceover:}  T ake this movement. I can stop on a frame that feels off, change it, and play the sequence again. Getting from ’something looks wrong’ to ’let’s try a fix’ feels much easier.\par

\textbf{Hình cần lấy:}  Chọn trước một chỉnh sửa thật: khoảng giữ pose hoặc một nét/điểm đặt chân. Xem trước → chỉnh → xem sau; dành ít nhất 3s cho bản sau.\par

\textbf{Dựng:}  Dùng nhãn Trước/Sau và cùng một vùng hình. Không gọi là cải thiện đã kiểm chứng nếu chỉ là thử nghiệm.\par

\textbf{Âm thanh:}  Hạ giọng/nhạc để 2–3s animation cuối được xem thoải mái.\par

\textbf{Chữ trên hình:}  One small adjustment\par

\Needspace{3\baselineskip}\textbf{S23 Từ Aseprite sang Unity}\par

05:15–05:30 · 15 giây · Tư liệu F13, F12\par

\textbf{English voiceover:}  Getting it into Unity became smoother in our workflow, too. I still check the canvas, the pivot, and where the feet sit. But seeing the updated animation in the game feels so satisfying.\par

\textbf{Hình cần lấy:}  Lưu file → đúng bước import đang dùng → nhân vật chạy trong Unity. Phóng điểm chân/pivot một lần, rồi tắt overlay.\par

\textbf{Dựng:}  Match cut cùng pose giữa hai phần mềm. Hiện đúng importer hoặc bước xuất sheet của project; không dựng kéo thả nếu thực tế còn bước khác.\par

\textbf{Âm thanh:}  Tiếng lưu file nối với click Play; giữ tiếng game ở cuối.\par

\textbf{Chữ trên hình:}  Canvas · pivot · in-game check\par

\textbf{Ghi chú:}  Aseprite là tên dự kiến thay cho “Aspire” trong bản nháp. Đối chiếu đúng phần mềm khi quay.\par

\clearpage

\Needspace{3\baselineskip}\textbf{S24 Bức hình đã chuyển động}\par

05:30–05:45 · 15 giây · Tư liệu F14, F15\par

\textbf{English voiceover:}  So now the character moves. We’re almost done, right? Well... a nice-looking animation and an attack that actually works in the game still need a few things to connect them.\par

\textbf{Hình cần lấy:}  Animation sạch 4s → cùng đòn đó trong editor → đứng ở frame sắp chạm địch. Bắt đầu hiện timeline logic.\par

\textbf{Dựng:}  Chuyển từ hình sạch sang có lớp giải thích ngay trên cùng khung, tránh mở chương bằng màn hình toàn code.\par

\textbf{Âm thanh:}  Một nhịp nghỉ trước “Well”; giữ nhạc tò mò.\par

\textbf{Chữ trên hình:}  Animation → attack logic\par

\Needspace{3\baselineskip}\textbf{9 Đẹp thôi chưa đủ: phải đánh được}\par

05:45–07:10 — Công cụ animation giúp tác giả kiểm soát đòn đánh.\par

\Needspace{3\baselineskip}\textbf{S25 Bắt đầu bằng công cụ có sẵn}\par

05:45–06:02 · 17 giây · Tư liệu F14\par

\textbf{English voiceover:}  I started with Unity’s built-in animation tools. They got me going. But for combat, I wanted to control exactly when things happened, in a way I could easily see and adjust.\par

\textbf{Hình cần lấy:}  Animator/Animation của bản cũ nếu còn → một động tác → dừng ở thời điểm cần gây sát thương. Code chỉ lướt 2s nếu giúp giải thích.\par

\textbf{Dựng:}  Không tô đỏ phần mềm như lỗi. Đánh dấu điều tác giả cần điều khiển, không liệt kê mọi menu.\par

\textbf{Âm thanh:}  Nhạc giữ nhịp, giọng giải thích chậm hơn ở ý chính.\par

\textbf{Chữ trên hình:}  Controlling the timing\par

\Needspace{3\baselineskip}\textbf{S26 Cho người không làm game hiểu}\par

06:02–06:19 · 17 giây · Tư liệu F15, F23\par

\textbf{English voiceover:}  For example, winding up shouldn’t deal damage yet. The hit area needs to appear during the swing, then the character needs time to recover. Small timing changes can make an attack feel very different.\par

\textbf{Hình cần lấy:}  Cùng một đòn ở ba đoạn: Wind-up → strike → recovery. Cho vùng đánh xuất hiện đúng đoạn giữa. Chạy lại một lần ở tốc độ thật.\par

\textbf{Dựng:}  Ba nhãn ngắn, một màu nhấn cho vùng đánh. Đây là sơ đồ giải thích; nếu chưa khớp build, ghi Minh họa.\par

\textbf{Âm thanh:}  Ba tiếng tick rất nhỏ; âm đòn thật ở lần phát tốc độ bình thường.\par

\textbf{Chữ trên hình:}  Wind-up → strike → recovery\par

\Needspace{3\baselineskip}\textbf{S27 Một công cụ hợp cách mình làm}\par

06:19–06:36 · 17 giây · Tư liệu F14\par

\textbf{English voiceover:}  I knew I could write more code for this. At the time, though, I wasn’t confident building the whole tool myself. Then I found an asset that suited our workflow and tried it out.\par

\textbf{Hình cần lấy:}  Tên asset thật/giao diện tool → nạp animation → chọn frame. Dành phần lớn hình cho công cụ đang hoạt động, không dừng lâu ở trang bán.\par

\textbf{Dựng:}  Cắt từ frame đã dừng ở cảnh trước sang cùng frame trong tool mới.\par

\textbf{Âm thanh:}  Nhạc có một lớp sáng hơn khi công cụ xuất hiện.\par

\textbf{Chữ trên hình:}  A tool that fits our workflow\par

\textbf{Ghi chú:}  Điền tên asset thật khi có footage; không đoán tên hoặc nói đây là công cụ tự viết toàn bộ.\par

\clearpage

\Needspace{3\baselineskip}\textbf{S28 Các hộp nhìn đơn giản mà rất hữu ích}\par

06:36–06:53 · 17 giây · Tư liệu F14, F15\par

\textbf{English voiceover:}  I added a few custom parts for our game. Now I can preview the animation and edit the collision areas right there: where an attack can hit, and where a character can be hit.\par

\textbf{Hình cần lấy:}  Chọn một frame → chỉnh hitbox → bật hurtbox → chạy preview. Hai vùng dùng màu và nhãn khác nhau, không chỉ phân biệt bằng màu.\par

\textbf{Dựng:}  Crop đúng khu vực làm việc, giấu bớt panel không liên quan. Giữ overlay vài giây đủ đọc.\par

\textbf{Âm thanh:}  Nhạc thấp, tiếng thao tác nhỏ. Ngừng lời 2s khi hai vùng cùng hiện.\par

\textbf{Chữ trên hình:}  Hitbox = attack area · Hurtbox = vulnerable area\par

\Needspace{3\baselineskip}\textbf{S29 Công cụ trả lại thời gian thử}\par

06:53–07:10 · 17 giây · Tư liệu F14, F01\par

\textbf{English voiceover:}  The best part is being able to experiment: adjust something, replay it, and see what’s still off. There’s still logic to write, but now I have a clearer way to work.\par

\textbf{Hình cần lấy:}  Thao tác chỉnh → preview → gameplay sạch 5–6s. Dùng một đòn xuyên suốt để người xem nhận ra kết quả.\par

\textbf{Dựng:}  Tắt lớp hitbox trên cùng khung rồi để game chạy. Chuyển chương bằng nhân vật chuẩn bị nhảy.\par

\textbf{Âm thanh:}  Tăng tiếng game trong 5s cuối, giọng dừng để khán giả quan sát.\par

\textbf{Chữ trên hình:}  Adjust → test → watch\par

\Needspace{3\baselineskip}\textbf{10 Nhảy lên dễ, đáp xuống mới khó}\par

07:10–08:05 — Một vấn đề vật lý cụ thể, có kết quả nhìn thấy.\par

\Needspace{3\baselineskip}\textbf{S30 Nhân vật đã đánh được nhưng còn phải đáp đất}\par

07:10–07:24 · 14 giây · Tư liệu F16, F22\par

\textbf{English voiceover:}  Then came physics. I thought making a character jump and fall back down would be fairly straightforward. Landing taught me otherwise. Sometimes, my character seemed to have its own interpretation of the rules.\par

\textbf{Hình cần lấy:}  Một lần nhảy bình thường → clip lỗi thật như trượt/nảy nếu có → phản ứng nhẹ của người kể.\par

\textbf{Dựng:}  Phát lỗi một lần đủ hiểu, tối đa replay thêm một lần. Nếu không có lỗi cũ, dùng bản test hiện tại với nhãn rõ.\par

\textbf{Âm thanh:}  Giữ tiếng tiếp đất; nhạc hạ ngắn ở lỗi.\par

\textbf{Chữ trên hình:}  The landing problem\par

\Needspace{3\baselineskip}\textbf{S31 Giải thích cảm giác thay vì dạy thông số}\par

07:24–07:38 · 14 giây · Tư liệu F16\par

\textbf{English voiceover:}  With my project’s scale and setup, I had to balance jump speed, gravity, and landing behavior. The goal was simple: the character should respond to your input and feel grounded when it lands.\par

\textbf{Hình cần lấy:}  Ba nhãn bên cạnh cùng một lần nhảy: lên → rơi → chạm đất. Inspector chỉ hiện phần đang chỉnh.\par

\textbf{Dựng:}  Không phủ con số 250 lên toàn màn hình. Có thể ghi nó như thông số thử nghiệm khi đã xác nhận đơn vị/cấu hình.\par

\textbf{Âm thanh:}  Giọng rõ, ít thuật ngữ; tiếng game nền nhỏ.\par

\textbf{Chữ trên hình:}  Jump · fall · land\par

\textbf{Ghi chú:}  Không dùng lập luận “1 Unity unit nên velocity phải 250”. Ma sát cũng không phải lời giải chung cho mọi lỗi tiếp đất.\par

\clearpage

\Needspace{3\baselineskip}\textbf{S32 Một đối chiếu có thể nhìn thấy}\par

07:38–07:52 · 14 giây · Tư liệu F16\par

\textbf{English voiceover:}  Here’s one test before and after an adjustment. I want the landing to feel cleaner, so you can move or attack again. A few seconds on screen can take a lot of tweaking.\par

\textbf{Hình cần lấy:}  Hai take cùng góc, cùng thao tác. Trước 5s → sau 5s → sau chạy liên tục thêm 4s. Chọn đúng thay đổi có thật.\par

\textbf{Dựng:}  Đánh dấu vùng chân và thời điểm chạm đất, tránh phân đôi nếu làm nhân vật quá nhỏ.\par

\textbf{Âm thanh:}  Lời dừng sớm, cho hai lần tiếp đất được nghe rõ.\par

\textbf{Chữ trên hình:}  Before the adjustment · current version\par

\textbf{Ghi chú:}  Nếu thiếu bản trước, thay lời cảnh này bằng: “Here’s the landing in the current build. I want it to feel clean, so you can move or attack again. A few seconds on screen can take a lot of tweaking.”\par

\Needspace{3\baselineskip}\textbf{S33 Khép câu đùa từ đầu}\par

07:52–08:05 · 13 giây · Tư liệu F16, F01\par

\textbf{English voiceover:}  So now we have a character that jumps... and mostly cooperates when landing. It sounds basic. But every time it works, I’m still happy, like I’ve finally solved a stubborn problem.\par

\textbf{Hình cần lấy:}  Một cú nhảy sạch, tiếp đất rồi bước/đánh ngay nếu build hỗ trợ. Kết ở tư thế bắt đầu tấn công.\par

\textbf{Dựng:}  Gọi lại hình ảnh nhảy ở mở đầu, lần này cho xem đầy đủ kết quả.\par

\textbf{Âm thanh:}  Tiếng tiếp đất nối sang tiếng kiếm ở chương sau.\par

\textbf{Chữ trên hình:}  Finally staying on the ground\par

\Needspace{3\baselineskip}\textbf{11 Cuối cùng, một đòn đánh ra hồn}\par

08:05–09:10 — Trả lại cho người xem gameplay và tiếng game.\par

\Needspace{3\baselineskip}\textbf{S34 Công thức nghe rất đơn giản}\par

08:05–08:20 · 15 giây · Tư liệu F15, F17\par

\textbf{English voiceover:}  Combat sounds simple: an attack area meets a target area, you register a hit, and the enemy reacts. Easy enough. Here’s what that looked like... before everything quite agreed with the plan.\par

\textbf{Hình cần lấy:}  Hai vùng chạm → clip test lỗi thật hoặc phản ứng lệch → cùng đòn ở bản hiện tại.\par

\textbf{Dựng:}  Chừa 1s trước lộ kết quả. Không dựng lỗi giả rồi gọi là lịch sử phát triển.\par

\textbf{Âm thanh:}  Hạ nhạc ở “Here’s what that looked like”; để tiếng game làm câu chốt.\par

\textbf{Chữ trên hình:}  Sounds simple enough\par

\Needspace{3\baselineskip}\textbf{S35 Cái khó nằm ở cảm giác và thời điểm}\par

08:20–08:35 · 15 giây · Tư liệu F15, F18\par

\textbf{English voiceover:}  I still have to check when a hit counts, whether the reaction feels right, and what happens at close range. An attack can work in the code and still feel wrong on screen.\par

\textbf{Hình cần lấy:}  Một take có overlay thời điểm va chạm → phản ứng đối thủ → bản gameplay sạch. Chỉ chọn các hành vi có thể chỉ ngay trên hình.\par

\textbf{Dựng:}  Dừng một lần ở frame chạm, sau đó phát tốc độ thật. Tránh thuyết trình dài trên một ảnh code.\par

\textbf{Âm thanh:}  Âm va chạm và phản ứng đúng clip. Không dùng dựng âm để giả tính năng chưa có.\par

\textbf{Chữ trên hình:}  Timing · clear reactions\par

\clearpage

\Needspace{3\baselineskip}\textbf{S36 Từ những phần rời thành một chuỗi}\par

08:35–08:50 · 15 giây · Tư liệu F01, F18\par

\textbf{English voiceover:}  But when the art, animation, hit areas, and reactions start working together, the character feels more alive. Here’s the current build. I’ll let it run for a bit so you can see for yourself.\par

\textbf{Hình cần lấy:}  2–3s nhắc nhanh sprite → tool → scene; sau đó bắt đầu một encounter thật với nhân vật ở vị trí dễ đọc.\par

\textbf{Dựng:}  Kết lời trước cảnh sau. Không cắt ngay giữa đòn để che kết quả.\par

\textbf{Âm thanh:}  Nhạc lùi xuống; ưu tiên tiếng game.\par

\textbf{Chữ trên hình:}  Development build\par

\Needspace{3\baselineskip}\textbf{S37 Để game tự nói}\par

08:50–09:10 · 20 giây · Tư liệu F01\par

\textbf{English voiceover:}  No narration. Let the gameplay and game audio play.\par

\textbf{Hình cần lấy:}  Một đoạn gameplay liền mạch 15–20s: tiếp cận → giao tranh → kết thúc ngắn hoặc đổi vị trí. Ưu tiên rõ hành động hơn nhiều hiệu ứng. Tắt debug nhưng giữ UI cần thiết.\par

\textbf{Dựng:}  Không speed ramp, không zoom giật, không montage nửa giây. Hình cuối giữ 1–2s để làm cầu sang phần người kể.\par

\textbf{Âm thanh:}  20s không lời. Tiếng game rõ; nhạc hòa nhẹ rồi dịu dần.\par

\textbf{Chữ trên hình:}  Divergency · work in progress\par

\textbf{Ghi chú:}  Nếu build chưa có encounter dài tốt, dùng hai đoạn 8–10s và cắt ở chỗ hành động đã kết thúc.\par

\Needspace{3\baselineskip}\textbf{12 Sáu năm có nghĩa là gì?}\par

09:10–10:35 — Chạm vào sự lo lắng, rồi nhìn lại bằng chứng tiến bộ.\par

\Needspace{3\baselineskip}\textbf{S38 Nhìn lại sáu năm}\par

09:10–09:27 · 17 giây · Tư liệu F09, F02\par

\textbf{English voiceover:}  This project has been with me for about six years. The ideas, artwork, game- play, and code have all changed. Putting old and new designs side by side reminds me how far we’ve come.\par

\textbf{Hình cần lấy:}  Ba cặp ảnh cũ/mới, mỗi cặp có cùng tư thế/kích thước tương đối. Dừng lâu nhất ở thay đổi rõ nhất.\par

\textbf{Dựng:}  Dùng so sánh cạnh nhau hoặc wipe nhẹ. Chỉ ghi năm khi có căn cứ.\par

\textbf{Âm thanh:}  Nhạc bớt nhịp, giữ giai điệu ấm. Lời không cần quá trầm.\par

\textbf{Chữ trên hình:}  The versions along the way\par

\Needspace{3\baselineskip}\textbf{S39 Áp lực nhìn người khác làm nhanh}\par

09:27–09:44 · 17 giây · Tư liệu F22, F06\par

\textbf{English voiceover:}  Meanwhile, game development keeps changing. I hear so much about AI and building things faster. Then I look at our project and wonder: after all this time, why is there still so much left?\par

\textbf{Hình cần lấy:}  Người kể trước bàn → project hiện tại → ghi chú việc còn làm. Nếu cần hình minh họa AI, chỉ dùng chữ trung tính, không lấy thumbnail người khác làm đối tượng so sánh.\par

\textbf{Dựng:}  Giữ máy 5–7s ở người kể; không dựng một loạt tiêu đề kích thích.\par

\textbf{Âm thanh:}  Nhạc thấp hơn rõ, vẫn còn hơi ấm.\par

\textbf{Chữ trên hình:}  Sometimes I feel too slow\par

\textbf{Ghi chú:}  Đây là cảm nhận cá nhân, không khẳng định AI làm xong game trong 24 giờ và không suy đoán chất lượng game của người khác.\par

\clearpage

\Needspace{3\baselineskip}\textbf{S40 Nói thật về phần khó}\par

09:44–10:01 · 17 giây · Tư liệu F22, F06\par

\textbf{English voiceover:}  Making a game is hard technically, but it’s hard mentally, too. Sometimes a small bug stops me for ages. Sometimes I don’t know where to begin. And sometimes, I struggle to recognize our progress.\par

\textbf{Hình cần lấy:}  Cận người kể hoặc khung bàn tĩnh. Xen đúng một cảnh ghi chú/chỉnh sửa, không cần montage tuyệt vọng.\par

\textbf{Dựng:}  Ít cắt hơn phần kỹ thuật. Giữ khoảng thở cuối câu, không chèn đùa.\par

\textbf{Âm thanh:}  Lời là chính; nhạc nền rất nhẹ. Không cố tạo cao trào buồn.\par

\textbf{Chữ trên hình:}  Không cần chữ thêm\par

\Needspace{3\baselineskip}\textbf{S41 Tự hào bằng những thứ đang có}\par

10:01–10:18 · 17 giây · Tư liệu F02, F01, F19\par

\textbf{English voiceover:}  But then I look at the old artwork, and watch those characters running and fighting now. We’ve made something real. There’s plenty left to do, and I can still be proud of what’s here.\par

\textbf{Hình cần lấy:}  Hình cũ đã xuất hiện đầu video → cùng nhân vật trong build → cảnh môi trường đẹp → một khoảnh khắc nhóm làm việc nếu có.\par

\textbf{Dựng:}  Trả lại đúng cặp cũ/mới của mở đầu để câu chuyện khép vòng.\par

\textbf{Âm thanh:}  Nhạc ấm dần; đưa tiếng game vào cuối câu.\par

\textbf{Chữ trên hình:}  Little by little, it took shape\par

\Needspace{3\baselineskip}\textbf{S42 Bước tiếp theo có hình hài}\par

10:18–10:35 · 17 giây · Tư liệu F01, F22\par

\textbf{English voiceover:}  Now I want to keep improving it, and eventually watch other people explore this world. There’s a long way to go. But I can open the game and say: this is what we’re making.\par

\textbf{Hình cần lấy:}  Người kể → đoạn chạy qua môi trường → tên game trong khung chơi. Nếu chưa có người chơi thật, không dựng cảnh phản ứng playtest giả.\par

\textbf{Dựng:}  Giữ một khung gameplay đủ lâu, dùng làm nền cho phần Kickstarter sắp vào.\par

\textbf{Âm thanh:}  Nhạc trở lại nhịp vừa, tích cực nhưng không kiểu quảng cáo dồn ép.\par

\textbf{Chữ trên hình:}  Still working on it\par

\clearpage

\Needspace{3\baselineskip}\textbf{13 Bước tiếp theo: Kickstarter}\par

10:35–11:30 — Nói rõ hỗ trợ giúp hoàn thiện điều gì và mời một hành động.\par

\Needspace{3\baselineskip}\textbf{S43 Vì sao cần Kickstarter}\par

10:35–10:55 · 20 giây · Tư liệu F01, F20\par

\textbf{English voiceover:}  We’re preparing for Kickstarter. That support would help us put more resources\par

\textbf{into the work ahead:}  refining combat, building stages, adding audio, and testing. We want to bring what you’ve just seen closer to a finished game.\par

\textbf{Hình cần lấy:}  Gameplay hiện tại → bốn hình tương ứng combat, màn chơi, âm thanh, test. Kết ở trang chiến dịch hoặc art đã chuẩn bị.\par

\textbf{Dựng:}  Nói mục đích trước khi hiện nút. Chỉ nêu các hạng mục khớp kế hoạch campaign khi xuất video.\par

\textbf{Âm thanh:}  Nhạc giữ năng lượng vừa, tiếng game nhỏ dưới giọng.\par

\textbf{Chữ trên hình:}  Combat · stages · audio · testing\par

\Needspace{3\baselineskip}\textbf{S44 Một lời mời rõ ràng}\par

10:55–11:15 · 20 giây · Tư liệu F20, F21\par

\textbf{English voiceover:}  If you like where Divergency is going, we’d love to have you with us. You can find our Kickstarter page in the description and sign up to be notified when it launches. Thanks for hearing our story.\par

\textbf{Hình cần lấy:}  Trang prelaunch thật → vị trí nút nhận thông báo → trở lại khung game và logo. Có thể giữ một đường dẫn ngắn nếu đã có.\par

\textbf{Dựng:}  Chỉ một hành động chính. Không xếp thêm lời mời subscribe, donate, Discord cùng lúc.\par

\textbf{Âm thanh:}  Giọng mời nhẹ, không tăng nhạc thành một đoạn bán hàng.\par

\textbf{Chữ trên hình:}  Follow our Kickstarter · link in the description\par

\textbf{Ghi chú:}  Câu này chỉ dùng khi trang prelaunch đã công khai. Có hai bản lời thay thế cho trạng thái khác ở phần cuối tài liệu.\par

\Needspace{3\baselineskip}\textbf{S45 Kết bằng một hành động nhỏ}\par

11:15–11:30 · 15 giây · Tư liệu F01, F06, F21\par

\textbf{English voiceover:}  All right, I’m getting back to work. See you in the next Divergency update.\par

\textbf{Hình cần lấy:}  Người kể quay về màn hình → click Play → nhân vật hoàn thành một đòn đánh hoặc tiếp đất sạch → logo và lời mời duy nhất giữ khoảng 5s.\par

\textbf{Dựng:}  Khép lại bằng kết quả vui từ câu đùa “stay on the ground”. Chỉ dùng một beat, đừng thêm after-credit kéo dài.\par

\textbf{Âm thanh:}  Câu nói khoảng 5s; phần còn lại tiếng game và nhạc kết mềm, để dư đuôi âm.\par

\textbf{Chữ trên hình:}  DIVERGENCY · See you in the next update\par

\Needspace{3\baselineskip}\textbf{14 Checklist thu thập tư liệu}\par

Ưu tiên 1 tạo bộ khung để dựng được video. Ưu tiên 2 làm câu chuyện sinh động hơn. Ưu tiên 3 là phần thêm, có thể bỏ hoặc thay.\par

\Needspace{3\baselineskip}\textbf{F01 Gameplay hiện tại}\par

Đã có tư liệu dùng được Ưu tiên 1 · Cần quay build Quay 3–5 take, mỗi take 20–40s. Có một chuỗi liền mạch 20s, một đòn rõ, một cảnh chạy qua môi trường. Giữ tiếng game riêng, tắt debug. Dùng ở S01, S03, S11, S17, S29, S33, S36, S37, S41, S42, S43, S45\par

\clearpage

\Needspace{3\baselineskip}\textbf{F02 Build và hình cũ}\par

Đã có tư liệu dùng được Ưu tiên 1 · Cần tìm tư liệu Tìm 3–5 clip/ảnh đại diện. Ghi ngày hoặc phiên bản nếu xác minh được; chọn một cặp cũ/mới dùng ở cả đầu và cuối. Dùng ở S01, S02, S05, S07, S10, S38, S41\par

\Needspace{3\baselineskip}\textbf{F03 LF2 nguồn cảm hứng}\par

Đã có tư liệu dùng được Ưu tiên 2 · Cần quay hoặc tìm clip của mình Tự quay khoảng 20–30s, lấy đoạn 5–12s dễ đọc. Ghi rõ tên LF2. Dùng ở S04\par

\Needspace{3\baselineskip}\textbf{F04 Trao đổi với Khoa}\par

Đã có tư liệu dùng được Ưu tiên 2 · Cần tìm hoặc dựng lại Chọn 2–3 câu chat thật liên quan hướng game, hoặc dựng lại hai bong bóng lời mới có nhãn. Chỉ đưa phần định chia sẻ công khai. Dùng ở S06, S11\par

\Needspace{3\baselineskip}\textbf{F05 Sửa lỗi Divergency trước/sau}\par

Đã có tư liệu dùng được Ưu tiên 2 · Quay cùng một cảnh trước và sau khi sửa lỗi Hai clip thật trước và sau khi sửa cùng một lỗi trong Divergency; giữ góc máy và tình huống giống nhau. Dùng ở S08\par

\Needspace{3\baselineskip}\textbf{F06 Cảnh bàn làm việc}\par

Đã có tư liệu dùng được Ưu tiên 1 · Cần quay mới\par

\textbf{6–8 clip, mỗi clip 8–15s:}  bật đèn, mở project, tay chuột/phím, ghi chú, bấm Play, góc sau vai. Quay vài góc trong một buổi. Dùng ở S02, S04, S08, S09, S10, S11, S12, S39, S40, S45\par

\Needspace{3\baselineskip}\textbf{F07 Nghiên cứu và mô phỏng}\par

Đã có tư liệu dùng được Ưu tiên 2 · Cần tìm hoặc quay Một màn hình nghiên cứu có thể công khai, một simulation liên quan nếu còn. Không cần phương trình phức tạp. Dùng ở S12, S13, S14\par

\Needspace{3\baselineskip}\textbf{F08 Concept và không khí}\par

Đã có tư liệu dùng được Ưu tiên 1 · Có file gợi ý để tuyển\par

\textbf{Tuyển 5–7 ảnh:}  toàn cảnh, một chi tiết môi trường, địch/boss. Ghi rõ concept art. Boss áo trắng là tùy chọn, bỏ được.\par

\clearpage

Dùng ở S06, S07, S14, S15, S16, S17, S18\par

\Needspace{3\baselineskip}\textbf{F09 Deep và Solei cũ mới}\par

Đã có tư liệu dùng được Ưu tiên 1 · Có file hiện tại gợi ý Hai chân dung hiện tại, ít nhất hai cặp trước/sau có thứ tự đúng. Bản cũ chưa tìm được thì không dựng thành lịch sử giả. Dùng ở S07, S19, S38\par

\Needspace{3\baselineskip}\textbf{F10 Một animation xuyên suốt}\par

Đã có tư liệu dùng được Ưu tiên 1 · Cần tuyển từ nguồn làm việc Chọn một nhân vật và một đòn làm ví dụ. Lưu pose, sprite sheet và clip animation sạch. Dùng ở S19, S20, S22\par

\Needspace{3\baselineskip}\textbf{F11 Workflow Photoshop}\par

Đã có tư liệu dùng được Ưu tiên 3 · Cần tìm Quay file thật nếu còn. Thiếu footage thì dùng sheet có sẵn và lời kể trải nghiệm, không quay loading giả. Dùng ở S20\par

\Needspace{3\baselineskip}\textbf{F12 Thao tác Aseprite}\par

Đã có tư liệu dùng được Ưu tiên 1 · Cần quay mới\par

\textbf{3 take, mỗi take 20–30s:}  mở timeline, sửa một chi tiết thật, play trước/sau. Con trỏ di chuyển chậm. Dùng ở S21, S22, S23\par

\Needspace{3\baselineskip}\textbf{F13 Đưa sprite vào Unity}\par

Đã có tư liệu dùng được Ưu tiên 1 · Cần quay mới Quay đầy đủ luồng lưu/import đang dùng và kiểm tra pivot/chân. Lấy cả đoạn kết chạy trong game. Dùng ở S23\par

\Needspace{3\baselineskip}\textbf{F14 Tool animation}\par

Đã có tư liệu dùng được Ưu tiên 1 · Cần quay mới Quay tên tool, chọn frame, chỉnh vùng đánh, preview. Ghi lại tên asset và phần nhóm tùy biến. Dùng ở S24, S25, S27, S28, S29\par

\Needspace{3\baselineskip}\textbf{F15 Hitbox và hurtbox}\par

Đã có tư liệu dùng được Ưu tiên 1 · Cần quay mới Quay cùng một đòn với overlay bật và tắt. Đủ thấy lúc vùng đánh bắt đầu/kết thúc; màu kèm nhãn chữ.\par

\clearpage

Dùng ở S24, S26, S28, S34, S35\par

\Needspace{3\baselineskip}\textbf{F16 Nhảy và tiếp đất}\par

Đã có tư liệu dùng được Ưu tiên 1 · Cần quay và tìm bản cũ Cùng một góc quay, 2–3 take hiện tại và clip trước khi sửa nếu có. Giữ cả 3s trước nhảy và sau tiếp đất. Dùng ở S03, S30, S31, S32, S33\par

\Needspace{3\baselineskip}\textbf{F17 Lỗi đánh nhau thật}\par

Đã có tư liệu dùng được Ưu tiên 2 · Cần tìm Tìm 1–2 lỗi có thể hiểu sau 2–4s. Không có thì dùng thử nghiệm chưa ổn với nhãn đúng, hoặc bỏ câu đùa lỗi. Dùng ở S34\par

\Needspace{3\baselineskip}\textbf{F18 Phản ứng khi trúng đòn}\par

Đã có tư liệu dùng được Ưu tiên 1 · Cần quay mới Quay một pha tiếp cận → ra đòn → phản ứng địch ở tốc độ thật. Chọn kết quả khớp lời thu. Dùng ở S35, S36\par

\Needspace{3\baselineskip}\textbf{F19 Dấu vết hai người cùng làm}\par

Đã có tư liệu dùng được Ưu tiên 3 · Tùy chọn Một ảnh/call ngắn hoặc tư liệu trao đổi đã chọn để công khai. Thiếu thì dùng art và ghi chú thay thế. Dùng ở S41\par

\Needspace{3\baselineskip}\textbf{F20 Kickstarter}\par

Đã có tư liệu dùng được Ưu tiên 1 · Cần chốt trước xuất Xác nhận trạng thái trang, link thật, nút hành động và các hạng mục sẽ gọi vốn. Quay trang chỉ khi dùng được tại ngày xuất bản. Dùng ở S43, S44\par

\Needspace{3\baselineskip}\textbf{F21 Logo Divergency}\par

Đã có tư liệu dùng được Ưu tiên 1 · Có file gợi ý để tuyển Một PNG nền trong độ phân giải đủ lớn, thêm một nền game sạch cho đoạn kết. Dùng ở S03, S07, S15, S44, S45\par

\clearpage

\Needspace{3\baselineskip}\textbf{F22 Giọng và hình người kể}\par

Đã có tư liệu dùng được Ưu tiên 1 · Cần thu và quay Thu lời nháp trước. Khi quay chính thức, ưu tiên mặt người ở mở, câu đùa sửa lỗi, đoạn sáu năm và kết. Nếu không lộ mặt, dùng lời thu + cảnh tay/bàn. Dùng ở S03, S08, S18, S30, S39, S40, S42\par

\Needspace{3\baselineskip}\textbf{F23 Đồ họa giải thích}\par

Đã có tư liệu dùng được Ưu tiên 2 · Cần dựng đơn giản Hai bong bóng chat, một đồ thị hội tụ/lệch, ba nhãn lấy đà/ra đòn/thu đòn. Chuyển động đơn giản, không cần dựng 3D. Dùng ở S13, S14, S15, S26\par

\Needspace{3\baselineskip}\textbf{15 Hướng dẫn dựng}\par

\textbf{Giọng kể:}  Lời thu bằng tiếng Anh. Dùng “I” cho trải nghiệm cá nhân và “we” cho công việc chung. Đọc như kể với một người bạn mê game, giữ khoảng thở và tiếng cười thật.\par

\textbf{Nhịp hình:}  Đổi hình theo ý, thường 4–7s với hình có thông tin mới. Cảnh cảm xúc có thể giữ 6–10s; cảnh đòn đánh giữ đủ đầu và cuối. Không đặt quy tắc cứ vài giây là phải zoom.\par

\textbf{Chuyển cảnh:}  Phần lớn là cắt thẳng. Ưu tiên nối cùng pose, cùng vị trí điểm sáng hoặc cùng động tác; để lời cảnh sau vào sớm 0,3–0,7s khi phù hợp, hoặc giữ tiếng cảnh trước qua điểm cắt. Đây là ý nghĩa của J-cut và L -cut.\par

\textbf{Điểm nối đáng làm:}  Nguồn cảm hứng → prototype Divergency cùng hướng di chuyển; bong bóng chat → concept; đồ thị → cảnh thế giới; artwork → pose trong Aseprite; Aseprite → Unity cùng pose; hitbox bật → gameplay sạch; ảnh cũ ở đầu → ảnh mới ở cuối.\par

\textbf{Hài:}  Tự trêu cách làm của mình, không chê tác giả LF2, Khoa hoặc người dùng AI. Chọn khoảng 4–5 nhịp vui nhẹ; mỗi câu đùa có một khoảng nghỉ rồi đi tiếp. Boss hở là cảnh tùy chọn.\par

\textbf{Nhạc và âm thanh:}  Ba vùng cảm xúc: mở có nhịp vui; art có không khí tối; cuối ấm và vững hơn. Nhạc nền lùi khi có lời, hạ mạnh khi phát tiếng game. Không cần whoosh cho mọi lần đổi hình.\par

\textbf{Chữ:}  Chữ thực sự xuất hiện trong video được đề xuất bằng tiếng Anh. Ghi chú “Không cần chữ thêm” chỉ là hướng dẫn dựng. Phụ đề tiếng Anh tối đa hai dòng, tránh che bàn chân, địch và UI đang giải thích.\par

\textbf{Khi quay màn hình:}  Ngang 1920 ×1080, chọn 60 fps nếu máy giữ ổn định; giọng thu riêng. Mỗi take dư 3–5s đầu và cuối. Với pixel art, tránh phóng lẻ tùy tiện làm hình nhòe.\par

\textbf{Khi thiếu tư liệu:}  Thiếu build cũ thì dùng ảnh cũ; thiếu chat thật thì dựng lại có nhãn; thiếu lỗi thì dùng test hiện tại; thiếu facecam thì dùng cảnh bàn. Concept và hoạt cảnh minh họa được ghi nhãn tương ứng.\par

\textbf{Thời lượng:}  Bản tiếng Anh có khoảng 1379 từ lời dẫn chính. 11 phút 30 giây là mục tiêu dựng; căn lại từng mốc sau khi thu nháp. Dành thời gian cho gameplay không lời, xem art và các khoảng nghỉ; không tăng tốc giọng để khớp mốc cũ.\par

\clearpage

\Needspace{3\baselineskip}\textbf{16 Thứ tự thu thập và quay}\par

Buổi 1 tuyển tư liệu (60–90 phút dự kiến): Gom F02, F04, F08, F09, F21. Chọn cặp cũ/mới cho mở và kết; chốt nhân vật/đòn dùng xuyên suốt. Buổi 2 quay gameplay (60–90 phút dự kiến): Quay F01, F15, F16, F18; tìm F17 nếu có. Xem lại ngay để chắc nhân vật, phản ứng và tiếp đất nhìn rõ. Buổi 3 quay công cụ (45–75 phút dự kiến): Quay F12, F13, F14 liên tiếp với cùng animation. Bổ sung Photoshop và simulation nếu có. Buổi 4 lời kể và người thật (45–75 phút dự kiến): Thu nháp để sửa câu vấp; sau đó thu chính từng chương, hai take mỗi đoạn khó. Quay F06/F22 cùng một bố trí đèn. Dựng thô trước (Theo tư liệu thực tế): Ghép giọng và hình chứng minh câu chuyện. Xem từ đầu tới cuối, cắt lặp ý; rồi mới thêm đồ họa, nhạc và phụ đề.\par

\textbf{Tên file gợi ý:}  F16\_nhay\_tiep\_dat\_sau\_take02.mp4. Giữ bản sạch chưa có caption và ghi build, ngày quay, điều clip cho thấy.\par

\textbf{Nếu cần gần 10 phút:}  Giảm 15s hồi tưởng LF2, 12s boss tùy chọn, 15s workflow Photoshop, 15s giới thiệu asset, 13s câu đùa tiếp đất và 20s lặp ý ở phần suy ngẫm. Viết nối lại câu quanh chỗ cắt; giữ đoạn gameplay liền mạch và mục đích Kickstarter.\par

\textbf{Nếu muốn gần 12 phút:}  Thêm 10s đọc chi tiết art, 10s đối chiếu animation, 10s gameplay. Chỉ thêm khi tư liệu cung cấp thông tin mới.\par

\Needspace{3\baselineskip}\textbf{17 Lời mời Kickstarter theo trạng thái trang}\par

Đang chuẩn bị chưa có trang công khai We’re preparing a Kickstarter campaign for Divergency. I’ll add the link to the description when it’s ready. If you’d like to follow our progress, subscribe for future updates. Thanks for being here.\par

\textbf{Chữ cuối:}  Follow the making of Divergency. Chưa có trang thì chưa hiện nút nhận thông báo Kickstarter. Trang nhận thông báo đã công khai If you like where Divergency is going, we’d love to have you with us. You can find our Kickstarter page in the description and sign up to be notified when it launches. Thanks for hearing our story. Đây là phương án dùng trong S44. Kiểm tra link trên điện thoại trước khi xuất. Chiến dịch đang mở Divergency is now live on Kickstarter. If you’d like to help us keep building the game, you’ll find our plans and reward tiers through the link in the description. Thanks for being part of this. Đổi câu đầu S43 thành “Our Kickstarter campaign is now live.” Chỉ chèn hạn kết thúc nếu đã có lịch chính thức.\par

\Needspace{3\baselineskip}\textbf{18 Chi tiết cần khớp trước khi xuất}\par

\textbf{Tên phần mềm:}  Bản này dùng Aseprite thay cho “Aspire”. Khi quay, xác nhận đúng ứng dụng bạn đang dùng. Aseprite hỗ trợ animation và sprite sheet. Nguồn chính thức Tên phần mềm\par

\clearpage

\textbf{Luồng vào Unity:}  Tài liệu Unity có 2D Aseprite Importer, các chế độ nhập và tùy chọn pivot. Video cần thể hiện đúng importer hoặc luồng xuất sprite sheet của project mình, không hứa mọi project đều tự động giống nhau. Nguồn chính thức Luồng vào Unity\par

\textbf{Tính năng animation:}  Lời dẫn nói về cách bạn muốn kiểm soát workflow, không khẳng định Unity không làm được logic theo thời điểm. Tên asset đã mua cần điền theo project.\par

\textbf{Vật lý:}  Bỏ suy luận rằng chọn tỉ lệ sprite là đủ để suy ra velocity 250 hoặc friction là cách sửa chung. Chỉ mô tả điều thấy được trong phép thử; nếu giữ số, ghi rõ cấu hình và đơn vị của bản đó.\par

\textbf{Liên tưởng tên game:}  Đồ thị là cách minh họa cảm hứng cá nhân. Không trình bày chaos và divergence như hai khái niệm đồng nghĩa.\par

\par

\textbf{Sáu năm và công việc nghiên cứu:}  Giữ “khoảng sáu năm”; thêm năm lên ảnh chỉ khi đúng lịch sử. Thu đúng tình trạng PhD của bạn tại ngày ghi hình.\par

\textbf{Kickstarter:}  Chọn một trong ba lời mời ở trên; kiểm tra link và hạng mục gọi vốn thực tế. Không thêm ngày mở, số tiền hay mốc release chưa được nhóm chốt.\par

\Needspace{3\baselineskip}\textbf{19 Điểm tìm tư liệu trong thư mục dự án}\par

Xem file trước khi chọn cho shot. Các đường dẫn dưới đây tính từ thư mục gốc dự án Divergency Reviewer Site. GIF hoặc ảnh minh họa không tự chứng minh tính năng đã chạy trong build.\par

\textbf{Deep:}  \path{imgs/chars/Deep/Deep.png}\par

\textbf{Animation Deep:}  \path{imgs/chars/Deep/deep_sword_1.gif}\par

\textbf{Animation Deep:}  \path{imgs/chars/Deep/deep_sword_2.gif}\par

\textbf{Solei:}  \path{imgs/UI/character_selected_story_sprites/character_art_solei_card.png}\par

\textbf{Bản đồ Deep và Solei :}  \path{imgs/Stage1/Team_Deep_Solei_Tutorial_Map.png}\par

\textbf{Hình đòn Deep :}  \path{imgs/Steam+kick/DEEP_run_attack_skill_combo_fix69_steam_1920x1080.png}\par

\textbf{Hình đòn Deep :}  \path{imgs/Steam+kick/DEEP_basicattack_LifeCut8_steam_1920x1080.png}\par

\textbf{Logo:}  \path{imgs/Steam+kick/library_assets/divergency_library_logo_1280x720.png}\par

\Needspace{3\baselineskip}\textbf{20 Tiêu đề và thumbnail}\par

• We Spent Six Years Making This Game\par

\textbf{• Divergency:}  Six Years of Making a Game After Work\par

• I Research Stability, Then Make a Game Called Divergency\par

\textbf{Thumbnail đề xuất:}  nhân vật cũ và mới đặt cạnh nhau, cùng kích thước; chữ “6 YEARS LATER”. Chỉ dùng cặp ảnh có thứ tự phát triển thật.\par

\clearpage

\Needspace{3\baselineskip}\textbf{21 Kiểm tra trước khi xuất}\par

Mở đầu có gameplay hiện tại và một tương phản cũ mới trong 12s đầu. Mỗi câu về kết quả có đúng hình/clip để người xem đối chiếu.\par

Tư liệu cũ, concept, minh họa và gameplay được gọi đúng loại. Lời dựng lại với Khoa không xuất hiện như tin nhắn thật. Có tối thiểu một đoạn 15–20s được xem gameplay không lời. Lời rõ trên loa điện thoại; nhạc/tiếng pop không làm giật mình. Không có đoạn giải thích kéo dài mà hình không cung cấp thêm điều gì. Câu cuối Kickstarter đúng trạng thái trang và chỉ có một hành động chính. Đã xem toàn bộ một lượt như người chưa biết Divergency.\par

\end{document}